\pagenumbering{roman}

\begin{abstract} 

\noindent \gls{ai} has evolved through repeated cycles of excitement and disappointment, but the recent rise of \gls{genai} and \gls{llms} has sparked unprecedented public interest.
This raises the question of whether the general public has a sufficient understanding of the capabilities and limitations of the technology, and if not, how we can help them get there. 

This thesis addresses the knowledge gap by presenting a taxonomy of misconceptions surrounding \gls{llms}, based on available literature. 
We also present a mediation tool that is designed to help users reach a more accurate understanding of the technology through an interactive learning experience where users can build their own \gls{llm} by selecting different features in a skill-tree, which are connected to scrollytelling chapters that explain the features in more detail. 
Finally, we present a study on the prevalence of misconceptions in the different dimensions of the taxonomy, and the effectiveness of the mediation tool in resolving these misconceptions.
The study consisted of 10 participants, who were interviewed before and after using the mediation tool. The interviews were coded based on the taxonomy, and the changes in misconceptions were analyzed.

The results show that a majority of the participants in the study had misconceptions such as anthropomorphization, real-time information access, and described use that could be categorized as cognitive passivity, though environmental impact appeared to be well understood.
The mediation tool was found to be effective in addressing misconceptions like anthropomorphization, hallucination blindness and communicating the tokenization process, but work can still be done in addressing training data and context window misconceptions.

\end{abstract}

\renewcommand{\abstractname}{Acknowledgements}
\begin{abstract}
	Looking back on my years at university, it is amusing to think that an early fascination with the Enigma machine first drew me towards mathematics and has now, years later, led me to write about the misconceptions surrounding systems whose public perception is still shaped by the imitation game.

	First and foremost, I would like to thank my supervisors, David and Robert, who met my curiosity with enthusiasm and gave me the freedom to explore the topic in my own way, while always providing guidance and support when needed.

	I would also like to thank my family and friends for their support in every aspect of life outside the writing of this thesis. The prolonged lunch and coffee breaks and the late nights at Integrerbar are among my fondest memories from my time at the university.

	Finally, I must acknowledge the LLMs for their existence, and for being the subject of this thesis. Sadly, they will never understand this thesis, and by the time it is submitted, they will already have forgotten it.
	
	\vspace{1cm}
	\hspace*{\fill}\texttt{Alexander}\\ 
	\hspace*{\fill}\today
\end{abstract}
\setcounter{page}{1}
\newpage